\documentclass[12pt,letterpaper, onecolumn]{exam}
\usepackage[utf8]{inputenc}
\usepackage[T1]{fontenc}
\usepackage{amsmath}
\usepackage{amssymb}
\usepackage[lmargin=71pt, tmargin=1.2in]{geometry}
\usepackage{xcolor}
\usepackage[brazilian]{babel}
\usepackage{natbib}
\usepackage{enumitem}
\usepackage{booktabs}
\usepackage[hidelinks]{hyperref}
\urlstyle{same}

\renewenvironment{solution}
{\par\noindent\textbf{R:}\enspace\color{blue}}
{\par\color{black}}

\thispagestyle{empty}

\begin{document}

\begingroup
\centering
\LARGE Lista 10 Gabarito - Econometria I - 2026\\[0.5em]
\large Prof: André Luiz Pereira Mancha \par
\large Monitor: Tuffy Licciardi Issa \par
\endgroup
\rule{\textwidth}{0.4pt}
\pointsdroppedatright
\printanswers
\renewcommand{\solutiontitle}{\noindent\textbf{Ans:}\enspace}
\newcommand{\qstar}{\hspace{-0.6em}\textsuperscript{*}\hspace{0.2em}}

\begin{questions}

\question (8.5, \citep{wooldridge2016}) A variável \textit{smokes} é uma variável binária igual a um se a pessoa fuma, e zero caso contrário. Utilizando os dados contidos no arquivo \texttt{SMOKE}, estimamos um modelo de probabilidade linear de \textit{smokes}:

\[
\widehat{smokes} = 0{,}656 - 0{,}069\,\log(cigpric) + 0{,}012\,\log(income) - 0{,}029\,educ
\]
\[
\hspace{1.7cm} (0{,}855)\hspace{1.0cm}(0{,}204)\hspace{1.0cm}(0{,}026)\hspace{1.0cm}(0{,}006)
\]
\[
\hspace{0.6cm} +\; 0{,}020\,age - 0{,}00026\,age^2 - 0{,}101\,restaurn - 0{,}026\,white
\]
\[
\hspace{1.7cm} (0{,}006)\hspace{1.0cm}(0{,}00006)\hspace{1.0cm}(0{,}039)\hspace{1.0cm}(0{,}052)
\]
\[
\hspace{1.7cm} [0{,}856]\hspace{1.0cm}[0{,}207]\hspace{1.0cm}[0{,}026]\hspace{1.0cm}[0{,}006]
\]
\[
\hspace{1.1cm} [0{,}005]\hspace{1.0cm}[0{,}00006]\hspace{1.0cm}[0{,}038]\hspace{1.0cm}[0{,}050]
\]
\[
n = 807,\quad R^2 = 0{,}062.
\]

A variável \textit{white} é igual a um se a pessoa for branca, e zero caso contrário; as outras variáveis independentes foram definidas no Exemplo 8.7. Tanto os erros padrão usuais quanto os robustos em relação à heteroscedasticidade estão informados.

\begin{enumerate}[label=(\roman*)]
    \item Existe alguma diferença importante entre os dois conjuntos de erros padrão?
    \item Mantendo os outros fatores fixos, se a educação aumentar em quatro anos, o que acontece com a probabilidade estimada de fumar?
    \item Em que ponto mais um ano de idade reduz a probabilidade de fumar?
    \item Interprete o coeficiente da variável binária \textit{restaurn} (uma variável \textit{dummy} igual a um se a pessoa viver em um Estado em que há restrições ao fumo em restaurantes).
    \item A pessoa número 206 no conjunto de dados tem as seguintes características: \(cigpric = 67{,}44\), \(income = 6.500\), \(educ = 16\), \(age = 77\), \(restaurn = 0\), \(white = 0\) e \(smokes = 0\). Calcule a probabilidade de fumar dessa pessoa e comente o resultado.
\end{enumerate}

\begin{solution}
\begin{enumerate}[label=(\roman*)]
    \item Não. Os erros padrão usuais e robustos são muito parecidos em praticamente todos os coeficientes. Portanto, a heteroscedasticidade não parece alterar muito a inferência nessa aplicação.

    \item Quatro anos adicionais de estudo alteram a probabilidade prevista em
    \[
    4(-0,029) = -0,116.
    \]
    Logo, a probabilidade estimada de fumar cai em \(0,116\), isto é, \(11,6\) pontos percentuais.

    \item O efeito marginal de \textit{age} é
    \[
    \frac{\partial \widehat{smokes}}{\partial age}=0,020-2(0,00026)age.
    \]
    Igualando a zero,
    \[
    age^\ast=\frac{0,020}{0,00052}\approx 38,46.
    \]
    Assim, acima de cerca de 38,5 anos, um ano adicional de idade reduz a probabilidade prevista de fumar.

    \item O coeficiente \(-0,101\) significa que viver em um Estado com restrições ao fumo em restaurantes está associado a uma probabilidade de fumar cerca de \(10,1\) pontos percentuais menor, mantidos os demais fatores constantes.

    \item Substituindo os valores:
    \[
    \widehat{smokes}
    =
    0,656 - 0,069\log(67,44) + 0,012\log(6500) - 0,029(16) + 0,020(77) - 0,00026(77^2).
    \]
    Isso fornece aproximadamente
    \[
    \widehat{smokes}\approx 0,005.
    \]
    A probabilidade prevista é muito baixa e está bastante próxima do valor observado \(smokes=0\). Além disso, o valor previsto está dentro do intervalo \([0,1]\), o que evita um dos problemas possíveis do modelo de probabilidade linear.
\end{enumerate}
\end{solution}

\vspace{0.7cm}

\question (8.7, \citep{wooldridge2016}) Considere um modelo em nível de empregado,
\[
y_{i,e} = \beta_0 + \beta_1 x_{i,e,1} + \beta_2 x_{i,e,2} + \cdots + \beta_k x_{i,e,k} + f_i + v_{i,e},
\]
em que a variável não observada \(f_i\) é um ``efeito da empresa'' para cada empregado da empresa \(i\). O termo de erro \(v_{i,e}\) é específico do empregado e da empresa. O erro de combinação é \(u_{i,e} = f_i + v_{i,e}\), como na equação (8.28).
\begin{enumerate}[label=(\roman*)]
    \item Suponha que \(\operatorname{Var}(f_i)=\sigma_f^2\), \(\operatorname{Var}(v_{i,e})=\sigma_v^2\), e que \(f_i\) e \(v_{i,e}\) sejam não correlacionados. Demonstre que \(\operatorname{Var}(u_{i,e})=\sigma_f^2+\sigma_v^2\). Chame esse valor de \(\sigma^2\).
    \item Agora suponha que, para \(e \neq g\), \(v_{i,e}\) e \(v_{i,g}\) sejam não correlacionados. Demonstre que \(\operatorname{Cov}(u_{i,e},u_{i,g})=\sigma_f^2\).
    \item Seja
    \[
    \bar{u}_i = m_i^{-1}\sum_{e=1}^{m_i} u_{i,e}
    \]
    a média dos erros de combinação em determinada empresa. Demonstre que
    \[
    \operatorname{Var}(\bar{u}_i)=\sigma_f^2+\frac{\sigma_v^2}{m_i}.
    \]
    \item Detalhe a relevância do item (iii) para a estimação por MQP usando dados nivelados pela média da empresa, em que o peso usado para a observação \(i\) é o tamanho da empresa.
\end{enumerate}

\begin{solution}
\begin{enumerate}[label=(\roman*)]
    \item Como \(u_{i,e}=f_i+v_{i,e}\),
    \[
    \operatorname{Var}(u_{i,e})
    =
    \operatorname{Var}(f_i)+\operatorname{Var}(v_{i,e})+2\operatorname{Cov}(f_i,v_{i,e}).
    \]
    Pela não correlação, \(\operatorname{Cov}(f_i,v_{i,e})=0\). Portanto,
    \[
    \operatorname{Var}(u_{i,e})=\sigma_f^2+\sigma_v^2.
    \]

    \item Para \(e\neq g\),
    \[
    \operatorname{Cov}(u_{i,e},u_{i,g})
    =
    \operatorname{Cov}(f_i+v_{i,e},f_i+v_{i,g}).
    \]
    Expandindo,
    \[
    \operatorname{Cov}(u_{i,e},u_{i,g})
    =
    \operatorname{Var}(f_i)+\operatorname{Cov}(f_i,v_{i,g})+\operatorname{Cov}(v_{i,e},f_i)+\operatorname{Cov}(v_{i,e},v_{i,g}).
    \]
    Todos os termos de covariância cruzada são nulos, restando
    \[
    \operatorname{Cov}(u_{i,e},u_{i,g})=\sigma_f^2.
    \]

    \item Temos
    \[
    \bar{u}_i = f_i + \frac{1}{m_i}\sum_{e=1}^{m_i} v_{i,e}.
    \]
    Assim,
    \[
    \operatorname{Var}(\bar{u}_i)
    =
    \operatorname{Var}(f_i)
    +
    \operatorname{Var}\left(\frac{1}{m_i}\sum_{e=1}^{m_i} v_{i,e}\right),
    \]
    pois \(f_i\) e os \(v_{i,e}\) são não correlacionados. Como os \(v_{i,e}\) são não correlacionados entre si,
    \[
    \operatorname{Var}\left(\frac{1}{m_i}\sum_{e=1}^{m_i} v_{i,e}\right)
    =
    \frac{1}{m_i^2}\sum_{e=1}^{m_i}\operatorname{Var}(v_{i,e})
    =
    \frac{m_i\sigma_v^2}{m_i^2}
    =
    \frac{\sigma_v^2}{m_i}.
    \]
    Logo,
    \[
    \operatorname{Var}(\bar{u}_i)=\sigma_f^2+\frac{\sigma_v^2}{m_i}.
    \]

    \item O item (iii) mostra que a variância do erro médio depende do tamanho da empresa. Empresas maiores geram médias mais precisas porque \(\sigma_v^2/m_i\) cai com \(m_i\). Portanto, ao nivelar os dados pela média da empresa, observações com maior \(m_i\) devem receber maior peso em MQP, pois possuem menor variância do erro.
\end{enumerate}
\end{solution}

\vspace{0.7cm}

\question (8.1, \citep{wooldridge2016}) Quais das seguintes alternativas são consequências da heteroscedasticidade?
\begin{enumerate}[label=(\roman*)]
    \item Os estimadores MQO, \(\hat{\beta}_j\), são inconsistentes.
    \item A estatística \(F\) usual não mais tem uma distribuição \(F\).
    \item Os estimadores MQO não são mais BLUE.
\end{enumerate}

\begin{solution}
As consequências corretas são \textbf{(ii)} e \textbf{(iii)}.

\begin{enumerate}[label=(\roman*)]
    \item \textbf{Falsa}. Com exogeneidade mantida, os estimadores MQO continuam consistentes.
    \item \textbf{Verdadeira}. Sob heteroscedasticidade, a estatística \(F\) usual deixa de ter a distribuição \(F\) exata.
    \item \textbf{Verdadeira}. MQO deixa de ser BLUE porque perde eficiência frente a estimadores lineares não viesados que exploram a estrutura da variância.
\end{enumerate}
\end{solution}

\vspace{0.7cm}

\question (8.3, \citep{wooldridge2016}) Verdadeiro ou falso? O método MQP é preferido ao MQO quando uma variável importante for omitida.

\begin{solution}
\textbf{Falso}. MQP é útil para lidar com heteroscedasticidade conhecida ou modelada, mas não corrige viés por variável omitida. Se a variável omitida estiver correlacionada com alguma regressora incluída, tanto MQO quanto MQP serão inconsistentes.
\end{solution}

\vspace{0.7cm}

\question\qstar (8.C13, \citep{wooldridge2016}) Use os dados do arquivo \texttt{FERTIL2} para esta questão.
\begin{enumerate}[label=(\roman*)]
    \item Estime o modelo
    \[
    children = \beta_0 + \beta_1 age + \beta_2 age^2 + \beta_3 educ + \beta_4 electric + \beta_5 urban + u
    \]
    e registre os erros padrão usuais e robustos em relação à heteroscedasticidade. Os erros padrão robustos são sempre maiores do que os não robustos?
    \item Adicione três variáveis \textit{dummy} de religião e teste se elas são conjuntamente significativas. Quais são os \(p\)-valores dos testes não robustos e robustos?
    \item A partir da regressão do item (ii), obtenha os valores ajustados \(\hat{y}_i\) e os resíduos \(\hat{u}_i\). Regrida \(\hat{u}_i^2\) sobre \(\hat{y}_i\) e \(\hat{y}_i^2\), e teste a significância conjunta dos dois regressores. Conclua que a heteroscedasticidade está presente na equação para \textit{children}.
    \item Você diria que a heteroscedasticidade encontrada no item (iii) é importante na prática?
\end{enumerate}

\end{questions}
\bibliographystyle{apalike}
\bibliography{refs}
\end{document}
